%\usepackage{color}
\subsection{Dokumentanalys: SFERA och gamification}\label{subsec:dokumentanalys}
För att fastställa SFERA-protokollets tekniska uppsättning har en sammanställning gjorts utifrån de verk från litteraturstudien som berörde detta.
För att ge en introduktion till gamification och hur det har tillämpats i DAS-applikationer har en sammanfattning gjorts.

\subsubsection{SFERA-kommunikation}
Kommunikationsprotokollet som används inom C-DAS kallas SFERA och är som nämnt tidigare i avsnitt~\ref{subsubsec:c-das-litteraturstudie} av UIC för att stärka järnvägsbranschens genom teknisk vägledning i hur datautbyte i detta sammanhang bör se ut~\cite{IRS90940}.
Syftet är att energieffektivisera tågkörning, förbättra punktlighet, och förenkla internationell järnvägstrafik~\cite{IRS90940}.
DAS-verktyg är en nyckelfaktor i effektiviseringprocessen som dessutom inte kostar stora summon pengar i investeringskostnader~\cite{IRS90940}.
Så mycket som 5--15\% energi kan sparas in med DAS-verktyg~\cite{IRS90940, Yang2013} vilket är påfallande då energikostnade kan uppgå till 10\% av ett järnvägsföretags totala kostnader~\cite{Cwil2018}.
Vidare spås C-DAS och SFERA bidra till förbättringar inom energibesparing, punktlighet, kapacitet, kraftförsörjning, trafikstörningshantering, trafiksäkerhet och koldioxidutsläpp~\cite{IRS90940}.
\subsubsection{SFERA:s avgränsning}
SFERA sträcker sig inte bortom datautbyte och avhandlar exempelvis inte presentation av data för lokförare, kvaliteten av hastighetsrekommendationer eller information för lokförare, lokförarens beteende eller handlingar, algoritmer för beräkning av körprofil och så vidare~\cite{IRS90940}.
Istället tar SFERA sikte på att definiera data längst bak i informationskedjan så att andra infrastrukturförvaltare och järnvägsföretag får så bra beslutsunderlag som möjligt~\cite{IRS90940}.
\subsubsection{SFERA:s datauppställning}
Mycket av kommunikationen mellan TMS och C-DAS-OB består av journey profile och status reports.
Det är detta som i praktiken skiljer S-DAS från C-DAS - det standardiserade kommunikationsprotokollet samt möjligheten att utöva det~\cite{IRS90940}.
Journey profile innehåller alla data som krävs för att C-DAS-OB ska kunna beräkna en bra körprofil~\cite{IRS90940}.
Status reports är kommunikation på andra håller där information om tågets framförande~\cite{IRS90940}.
En annan typ av kommunikation är related train information, som delges C-DAS-OB från IM DAS-TS om innehåller positioner och hastigheter för kringliggande tåg på samma bana~\cite{IRS90940}.
C-DAS-OB kan då använda denna information för att räkna ut ny körprofil~\cite{IRS90940}.
Informationen kan även presenteras för föraren för att stärka dennes context awareness.
Vidare är en följd av related train information good-will hos förarna~\cite{IRS90940} som har visat starkt intresse för denna typ av data.
Datameddelanden inom SFERA är överlag designade att vara så små som möjligt för att öka flexibiliteten~\cite{IRS90940}.
\paragraph{}
\hspace{-11pt}En journey profile (se figur~\ref{tab:jp}) består av flera delar.
Tågfärdens unika referens-ID\footnote{Ett för det innevarande trafikdygnet unikt nummer mellan en och sex siffror.}~och information om tågets fysiska egenskaper (train characteristics, TC).
Se beskrivning av TC i figur(se figur~\ref{tab:tc}).
Den innehåller även statisk och dynamisk data om tågfärdens innehåll.
Statisk data benämns segment profile (se figur~\ref{tab:sp}) och är bland annat trafikplatser, lutningsförhållanden, kurvradier, ID för aktuella bana/banor, hastighetsbegränsningar, referenspunkter för C-DAS-OB:s positioneringssystem), signal- och balisplaceringar och mycket mer.
Dynamisk data är inte en egen datastruktur men kan röra sig om information tillfälliga hastighetsnedsättningar eller dåliga adhesionsförhållanden inom ett visst område.
Vidare innehåller JP tidskorridoren ett tåg måste hålla sig inom för att upprätthålla sin rättidighet.
\begin{table}[hbt!]
    \centering
    \begin{tabular}{|l|}
        \hline
        \textbf{Journey profile} \\
        \hline
        - Unikt ID för berörd tågfärd \\
        - Tidskorridor \\
        - Train characteristics \\
        - Segment profile (lista av) \\
        - Dynamisk data \\
        \hline
    \end{tabular}
    \caption{Datastruktur, journey profile}
    \label{tab:jp}
\end{table}
Status report (se figur~\ref{tab:sr}) är sättet på vilket C-DAS-OB kan skicka information till TMS\cite{IRS90940}.
Det består av tågets position i förhållande till referenspunkter från tilldelad SP (som del av JP), innevarande hastighet, senast passerade trafikplats och prognos för ankomst till nästa trafikplats, aktuell JP, förväntad energikonsumtion och -återmatning (elbroms) för kommande minuter samt statusmeddelanden\cite{IRS90940}.
Lokförare kan även rapportera in om förändringar i adhesionsförhållanden~\cite{IRS90940}.
Detta kommer skickas till TSM som sedan delger andra C-DAS-OB om detta via related train information.
Statusmeddelanden kan användas som notis om klart för avgång, om tågfärden avbryts och förseningsorsaker.
\begin{table}
    \centering
    \begin{tabular}{|l|}
        \hline
        \textbf{Status report} \\
        \hline
        - Position \\
        - Hastighet \\
        - ID för aktuell JP \\
        - Ändringar av TC \\
        - Förväntad energikonsumtion och -återmatning \\
        - Förändrade adhesionsförhållanden \\
        - Statusmeddelanden \\
        \hline
    \end{tabular}
    \caption{Datastruktur, status report}
    \label{tab:sr}
\end{table}
\begin{table}
    \centering
    \begin{tabular}{|l|l|l|}
        \hline
        \textbf{Segment profile} & & \\
        \hline
        - Unikt ID & & \\
        - Version & & \\
        - Längd & & \\
        \hline
        SP-punkter & SP-karaktäristik & SP-områden \\
        \hline
        - Trafikplats & - Banans STH för ett visst segment & - Station, plattform \\
        - Uppehåll\footnote{Plats för exempelvis resandeutbyte (ja/nej) eller tågmöte.} & - Unikt ID för aktuell bana & - Tunnel \\
        - Signalplacering & - Kurvradie, lutning & - Strömlös sektion \\
        - Balisplacering & - Kontaktledningsspänning & - Tillåten axellast \\
        - Referenspunkt (kilometer) & - Infrastrukturbegränsningar &  \\
        - Oövervakad plankorsning & - Tillgänglig(a) tågskyddssystem & \\
        \hline
    \end{tabular}
    \caption{Datastruktur, segment profile}
    \label{tab:sp}
\end{table}
\begin{table}
    \centering
    \begin{tabular}{|l|}
        \hline
        \textbf{Train characteristics} \\
        \hline
        Obligatoriskt: \\
        - Unikt ID \\
        - Version \\
        - Järnvägsföretags unika ID \\
        \hline
        Valfritt: \\
        - Typ av tågfärd\footnote{Persontåg, godståg, tjänstetåg, underhållsfordon} \\
        - STH \\
        - Motoreffekt (kW) \\
        - Typ av tågskyddssystem\footnote{ETCS, ATC (SE), ATC (NO). ATC (DK), osv.} \\
        - Tågvikt \\
        - Dynamisk tågvikt \\
        - Tåglängd \\
        - Maximal accelerationspotential \\
        - Komfortmässig accelerationspotential \\
        - Bromskraft \\
        - Tillsättningstid för bromssystem \\
        \hline
    \end{tabular}
    \caption{Datastruktur, train characteristics}
    \label{tab:tc}
\end{table}
\subsubsection{Gamification och DAS}
\paragraph{}
\hspace{-11pt}
Detta avsnitt är kopplat till forskningsfråga 4 som lyder \emph{\lq{}Hur kan gamification användas för att stödja och motivera lokförare i deras användande av C-DAS-verktyg?\rq{}}.
Det finns exempel på när gamification tillämpats i DAS-applikationer, men de är få till antalet~\cite{Bae2014, Cwil2018}.
Två studier har berörda två olika digitala lokförarverktyg som tagits fram.
Den ena är en applikation för ett stort polskt regionalt järnvägsföretag implementerade ett eco driving\footnote{Energieffektiv körstil som även sparar in på fordons- och infrastrukturslitage.}-verktyg i sin verksamhet 2018~\cite{Cwil2018}, och den andra en prototyp som har tagits fram i Norge för att främja lokförarnas situationsmedvetenhet~\cite{Bae2014}.
\paragraph{}
\hspace{-11pt}För den polska applikationen introducerades poängbaserade gamification-element stegvis i fyra olika faser med 2--3 månaders mellanrum så att förarna inte skulle bli överväldigare och lära sig systemet i en god takt.
Syftet var som tidigare nämnt att förbättra förarnas eco driving.
Alla faser har haft internalisering som mål och har därför behövt tillgodose förarnas behov av autonomi, kompetensprövning och känsla av att vara i ett sammanhang (samhörighet).
För samtliga faser har användandet av verktyget varit frivilligt för att säkerställa täckning för autonomibehovet.
\paragraph{}
\hspace{-11pt}Den första fasen syftade till att synliggöra för förarna hur energieffektivt de körde genom att mäta hur mycket effekt från kontaktledningen som använts och återge detta så snabbt som möjligt.
Meningen var inte att det skulle jämföras mot andra förare eller uppnå ett särskilt mål, utan snarare att skapa en \lq{}placebo-effekt\rq{} genom att göra dem uppmärksamma på att deras effektivitet mättes och sparades.
Enkom genom detta sjönk energiförbrukningen med 8\% på en av linjerna där pilottester kördes~\cite{Cwil2018}.
Att den data som presenterades var korrekt visade sig vara avgörande för lokförarnas förtroende för systemet; engagemanget sjönk avsevärt när förarna inte trodde att deras angivna förbrukning inte stämde överens med verkligheten.
Denna gamification-komponent är i linje med vad SDT säger om autonomi och motivation då deltagande var frivilligt och inga negativa konsekvenser följde av att inte deltaga.
Belöningar kopplade till användning tillämpades för att stimulera kompetensprövning.
Känslan av samhörighet influerades genom att informera om att all förare fanns med i systemet, även om deltagandet var frivilligt och information om andra förares prestationer inte fanns tillgängligt.
\paragraph{}
\hspace{-11pt}Andra fasen syftade till att introducera gamification-moment där förarna kunde jämföra sig med varandra.
Individuella prestationer publicerades för alla att se och uppdrag och utmaningar gjordes tillgängliga.
Typiska diton kunde vara att spara 100~kWh på en vecka, förbättra sin genomsnittliga energikonsumtion eller minska konsumtionen på en viss bana med 5\%~\cite{Cwil2018}.
Belöningar för de högst presterande förarna under en given period gavs och istället för att belönas för användande eller prestation gavs de istället för att köra en särskild bana på effektivast vis.
Deltagande var fortfarande frivilligt, och val av utmaningar eller uppdrag likaså för att hedra förarnas autonomi.
Samhörigheten närmades genom att visa individuell statistik jämfört med genomsnittsföraren samt publicera de bäst presterande förarna och deras statistik.
Kompetensprövning stimulerades genom känslan av progression när uppdrag och utmaningar klarades av och blev svårare.
\paragraph{}
\hspace{-11pt}I tredje fasen låg fokus på att stärka förarnas intrinsiska motivation genom laganda och samarbete.
Förarna tävlade tillsammans i grupper som utsågs av instruktörer för att nå specifika lagorienterade mål såsom att spara 10~000~kWh på en månad.
Tävlingsmoment och kompetensprövning riktades istället mot andra grupper samt att gruppernas poängställning offentliggjordes på en topplista.
Genom att gå från individuella tävlingar till grupptävlingar tillgodosågs samhörighetsbehovet.
\paragraph{}
\hspace{-11pt}För fjärde och sista fasen låg fokus på att aktivera de bästa förarna på ett sätt så att de skulle få ett fortsatt intresse av att fortsätta använda systemet.
Lösningen var att dem att hjälpa andra förare att bli bättre på att använda verktygen och i förlängningen eco driving och stimulera båda parters samhörighetskänsla.
Istället för att kompetensprövningen låg direkt i systemet skedde det indirekt genom kunskapsöverföring utanför verktyget.
De kunde även skänka \lq{}altruismpoäng\rq{} till andra förare för att behålla deras incitament att fortsätta att tjäna in poäng.
\paragraph{}
\hspace{-11pt}För den norska prototypen var syftet att stärka lokförarnas situationsmedvetenhet.
Ett poängsystem nämns inte uttryckligen men är underförstått som ett sätt att mäta prestationer i form av exempelvis punktlighet eller effektiv körstil.
För detta implementerades en systemkomponent som visuellt kunde återge visualisera trafiken i tågets närområde.
Verktygets kärna låg i jämförelsen med kollegor där förarna kunde se hur deras handlingar påverkat deras punktlighet och prestation, samt hur andra förare tenderat att agera i liknande situationer.
Reflektion var en viktig del där verktyget stöttade individen i att utvärdera den egna körstilen i en strävan att förbättra den.
Situationer kunde då spelas upp i efterhand för analys av beslut och eventuella förbättringar.
Systemet skulle vara engagerande och interaktivt för att stimulera intresset och uppmärksamheten av verktyget utan att störa förarens uppmärksamhet.
\paragraph{}
\hspace{-11pt}Båda projekten har olika syften och utvecklades oberoende av varandra men visar på gemensamma lärdomar och strategier.
Exempelvis har båda systemen tagits fram med gamification som primärt fokus få att uppnå sitt mål, vare sig det varit eco driving eller ökad sitautionsmedvetenhet~\cite{Bae2014, Cwil2018}.
Båda studierna nämner vikten av feedback, och att förarna får god möjlighet att utvärdera sig själv på ett autonomt sätt~\cite{Bae2014, Cwil2018}.
Att gamification-elementen är sömlöst integrerade och inte stör lokförarens dagliga arbete, blir ett irritationsmoment eller skapar onödig distraktion har i båda fallen visat sig vara av stor vikt~\cite{Bae2014, Cwil2018}.
Vidare har tillämpningen i båda projekten delvis vilat på beteendevetenskapliga teorier om motivation och hur den bäst främjas för långvarighet och att lokförarna inte ska \lq{}tröttna\rq{} på verktyget efterhand~\cite{Bae2014, Cwil2018}.
I båda studier framhålls vikten av att tillgodose tre mänskliga behov enligt självbestämmandeteorin för att säkerställa den bästa formen av motivation; autonomi, känslan av samhörighet samt kompetensprövning.
Slutsatserna från båda studier visar två gemensamma punkter.
Gamification ~\emph{kan} användas i digitala verktyg för lokförare för att effektivt motivera dem och öka deras prestation~\cite{Bae2014, Cwil2018}.
Faktumet att det visat sig vara effektivt tyder på att det finns fog att undersöka ämnet vidare och att det finns potential för bredare tillämpning i järnvägsbranschen~\cite{Bae2014, Cwil2018}.
I båda studier konstateras även att gamification i DAS-system för lokförare är svårt.
Det fordrar mycket god domänkunskap, tydliga och noggrant avvägda begränsningar i komponenter för att anpassa systemet efter verksamheten samt vidmakthållandet av förarnas förtroende~\cite{Bae2014, Cwil2018}.
%Jobbar inte med datorer, lokförare
%Tågläge styrker och legitimerar körprofil